\documentclass[10.5pt]{article}
\usepackage[letterpaper, margin=0.6in, top=0.5in, bottom=0.5in]{geometry}
\usepackage{amsmath,amssymb,amsthm}
\usepackage{fancyhdr}
\usepackage{xcolor}
\usepackage{tikz}
\usepackage{enumitem}
\usepackage{tcolorbox}
\usepackage{hyperref}

% Define colors matching instructor board annotations
\definecolor{chalkblue}{RGB}{0, 105, 185}
\definecolor{chalkred}{RGB}{205, 35, 35}
\definecolor{chalkpink}{RGB}{230, 90, 120}
\definecolor{chalkcyan}{RGB}{15, 160, 190}
\definecolor{chalkgreen}{RGB}{20, 130, 50}
\definecolor{boxbg}{RGB}{252, 252, 254}
\definecolor{boxframe}{RGB}{150, 160, 180}

\pagestyle{fancy}
\fancyhf{}
\renewcommand{\headrulewidth}{0pt}
\renewcommand{\footrulewidth}{0pt}

\setlength{\parindent}{0pt}
\setlength{\parskip}{2.5pt}

\begin{document}

% ==============================================================================
% PAGE 1
% ==============================================================================
\noindent \textbf{\large Math 105L} \hfill \textbf{\large Lecture 5-1}\\[1pt]
{\textbf{\Large Limits, Continuity, \& One-Sided Limits}}\\[-5pt]
\rule{\textwidth}{1.2pt}

\vspace{0.03in}

\textbf{\large 1. The Three Conditions for Continuity at $x = c$}

A function $f(x)$ is continuous at a point $x = c$ if and only if:
\begin{enumerate}[leftmargin=1.4em, topsep=1pt, itemsep=2pt]
  \item $\mathbf{f(c)}$ is defined (the function exists at the point).
  \item $\mathbf{\lim_{x\to c} f(x)}$ exists (the values approach a single finite real number).
  \item $\mathbf{\lim_{x\to c} f(x) = f(c)}$ (what we expect the function to be matches what it actually is).
\end{enumerate}

\vspace{0.03in}

\begin{tcolorbox}[colback=boxbg, colframe=boxframe, boxrule=0.8pt, arc=3pt, left=8pt, right=8pt, top=4pt, bottom=4pt]
\textbf{\large Three Ways Continuity Can Fail (Types of Discontinuities)}
\begin{itemize}[leftmargin=1.2em, topsep=2pt, itemsep=4pt]
  \item \textbf{1. Hole / Removable Discontinuity:}
    We know what to expect as $x \to c$, but there is a hole in the graph ($f(c)$ is undefined or detached from the curve).
  \item \textbf{2. Infinite Discontinuity (Vertical Asymptote):}
    The values grow without bound ($f(x) \to \pm\infty$ as $x \to c$).
  \item \textbf{3. Jump Discontinuity:}
    The function approaches two completely different finite values from the left and from the right ($\lim_{x\to c^-} f(x) \ne \lim_{x\to c^+} f(x)$).
\end{itemize}
\end{tcolorbox}

\vspace{0.04in}

\textbf{\large 2. One-Sided Limits \& Two-Sided Limit Existence}

\begin{tcolorbox}[colback=chalkblue!6, colframe=chalkblue, boxrule=1pt, arc=3pt, left=8pt, right=8pt, top=4pt, bottom=4pt]
\textbf{The Fundamental Criterion for Limit Existence:}\\[2pt]
$\lim_{x\to c} f(x) = L$ exists if and only if both one-sided limits exist and are equal to $L$:
\[
  \mathbf{\lim_{x\to c} f(x) = L \iff \lim_{x\to c^+} f(x) = L = \lim_{x\to c^-} f(x)}
\]
\begin{itemize}[leftmargin=1.2em, topsep=1pt, itemsep=2pt]
  \item $\mathbf{\lim_{x\to c^+} f(x)}$: Right-hand limit (approaching $c$ from numbers greater than $c$).
  \item $\mathbf{\lim_{x\to c^-} f(x)}$: Left-hand limit (approaching $c$ from numbers less than $c$).
\end{itemize}
\end{tcolorbox}

\vspace{0.04in}

\textbf{\large 3. Worked Classroom Analysis: Piecewise \& Rational Limits}

\begin{itemize}[leftmargin=1.2em, topsep=1pt, itemsep=4pt]
  \item \textbf{Example 1: The Signum Function $f(x) = \frac{|x|}{x}$ at $x = 0$:}
  Using the piecewise definition $|x| = \begin{cases} x, & x \ge 0 \\ -x, & x < 0 \end{cases}$:
  \[
    f(x) = \begin{cases} \frac{x}{x} = 1, & x > 0 \\ \frac{-x}{x} = -1, & x < 0 \end{cases}
  \]
  \begin{itemize}[label=--]
    \item Right-hand limit: $\lim_{x\to 0^+} \frac{|x|}{x} = \lim_{x\to 0^+} (1) = \mathbf{1}$.
    \item Left-hand limit: $\lim_{x\to 0^-} \frac{|x|}{x} = \lim_{x\to 0^-} (-1) = \mathbf{-1}$.
    \item Since $1 \ne -1$, the two-sided limit $\mathbf{\lim_{x\to 0}\frac{|x|}{x} \text{ Does Not Exist (DNE)}}$ (Jump Discontinuity).
  \end{itemize}

  \item \textbf{Example 2: The Reciprocal Function $f(x) = \frac{1}{x}$ at $x = 0$:}
  \begin{itemize}[label=--]
    \item As $x \to 0^+$, denominators are positive and shrink to 0: $\lim_{x\to 0^+} \frac{1}{x} = \mathbf{+\infty}$.
    \item As $x \to 0^-$, denominators are negative and shrink to 0: $\lim_{x\to 0^-} \frac{1}{x} = \mathbf{-\infty}$.
    \item The two-sided limit $\mathbf{\lim_{x\to 0}\frac{1}{x} \text{ DNE}}$ (Infinite Discontinuity / Vertical Asymptote).
  \end{itemize}
\end{itemize}

\end{document}

